\subsection{Mathematical Foundations and Structural Syntax}

A \texttt{set} represents a domain of unique elements. The word ``domain'' is useful here because a set is not mainly about position, repetition, or sequence traversal. It is about whether a value belongs to a collection.

A list may contain the same value many times. A set does not. If the same value is inserted repeatedly, the set keeps only one logical occurrence.

\begin{quote}
A list remembers positions and repetitions. A set remembers membership.
\end{quote}

This makes sets useful for questions such as:

\begin{itemize}
    \item has this value already appeared?
    \item which unique values exist in this input?
    \item which values belong to both groups?
    \item which values belong to one group but not another?
\end{itemize}

The mathematical idea is old, but Python implements it with a runtime hash table. That means a \texttt{set} is not merely a mathematical notation. It is also a concrete memory structure with hash-based placement rules.

\subsubsection{Defining Unique Unordered Domains Using the Curly-Brace \texttt{\{\}} Construct}

A non-empty set can be written with curly braces. The elements are separated by commas.

\begin{verbatim}
facts = {"spam", "eggs", "spam", 42}

print(f"facts: {facts!r}")
print(f"type(facts): {type(facts)}")
print(f"len(facts): {len(facts)}")
\end{verbatim}

Possible output:

\begin{verbatim}
facts: {'eggs', 42, 'spam'}
type(facts): <class 'set'>
len(facts): 3
\end{verbatim}

The input contained \texttt{"spam"} twice, but the resulting set contains it only once. The duplicate was not stored as a second independent component because sets enforce uniqueness.

The printed order should not be interpreted as a stable structural order. A set is not a list with unusual brackets. It is not a tuple with mutable syntax. It is a hash-based uniqueness domain.

We can inspect the object with \texttt{dir()}.

\begin{verbatim}
facts = {"spam", "eggs", 42}

print(f"type(facts): {type(facts)}")
print()

print("__contains__" in dir(facts))
print("__iter__" in dir(facts))
print("__len__" in dir(facts))
print("add" in dir(facts))
print("remove" in dir(facts))
print("union" in dir(facts))
print("intersection" in dir(facts))
\end{verbatim}

Possible output:

\begin{verbatim}
type(facts): <class 'set'>

True
True
True
True
True
True
True
\end{verbatim}

The most important names at this stage are:

\begin{itemize}
    \item \texttt{\_\_contains\_\_}: supports membership testing with \texttt{in};
    \item \texttt{\_\_iter\_\_}: allows traversal over the current set elements;
    \item \texttt{\_\_len\_\_}: supports \texttt{len()};
    \item \texttt{add} and \texttt{remove}: later used for direct mutation;
    \item \texttt{union} and \texttt{intersection}: later used for mathematical set calculations.
\end{itemize}

Membership is the natural first operation for a set.

\begin{verbatim}
facts = {"spam", "eggs", 42}

print(f"{'spam' in facts = }")
print(f"{'larch' in facts = }")
print(f"{42 in facts = }")
\end{verbatim}

Possible output:

\begin{verbatim}
'spam' in facts = True
'larch' in facts = False
42 in facts = True
\end{verbatim}

This code does not ask where \texttt{"spam"} is located. It asks whether \texttt{"spam"} belongs to the set. That is the correct mental model.

A set may contain different hashable object types.

\begin{verbatim}
mixed = {"D'oh!", 42, ("Arthur", "Brian")}

print(f"mixed: {mixed!r}")
print(f"type(mixed): {type(mixed)}")

for item in mixed:
    print(f"item={item!r:<22} type={type(item)}")
\end{verbatim}

Possible output:

\begin{verbatim}
mixed: {42, "D'oh!", ('Arthur', 'Brian')}
type(mixed): <class 'set'>
item=42                     type=<class 'int'>
item="D'oh!"                type=<class 'str'>
item=('Arthur', 'Brian')    type=<class 'tuple'>
\end{verbatim}

The tuple is accepted because it is immutable and its contents are hashable. A list would not be accepted.

\begin{verbatim}
try:
    bad_set = {"spam", ["eggs", 42]}
except TypeError as err:
    print(f"set construction failed: {err}")
\end{verbatim}

Possible output:

\begin{verbatim}
set construction failed: unhashable type: 'list'
\end{verbatim}

The list is rejected because a list can mutate. A mutable object cannot safely act as a hash-table element when its hash-relevant state may change after insertion.

\subsubsection{Instantiation Boundaries: Differentiating Empty Set Initialization \texttt{set()} from Empty Dictionary Literal Declaration \texttt{\{\}}}

The syntax for a non-empty set uses curly braces, but the syntax for an empty set does not.

This is a deliberate boundary in Python syntax:

\begin{itemize}
    \item \texttt{\{\}} creates an empty dictionary;
    \item \texttt{set()} creates an empty set.
\end{itemize}

The distinction exists because curly braces were already used for dictionaries before set literals were added to Python. Therefore, the empty curly-brace form kept its dictionary meaning.

\begin{verbatim}
empty_a = {}
empty_b = set()

print(f"empty_a: {empty_a!r}")
print(f"type(empty_a): {type(empty_a)}")
print()

print(f"empty_b: {empty_b!r}")
print(f"type(empty_b): {type(empty_b)}")
\end{verbatim}

Possible output:

\begin{verbatim}
empty_a: {}
type(empty_a): <class 'dict'>

empty_b: set()
type(empty_b): <class 'set'>
\end{verbatim}

This is one of the common beginner traps. The expression \texttt{\{\}} does not mean ``empty curly-brace collection in general.'' It specifically means empty dictionary.

Because this subsection introduces both objects syntactically, we can inspect both types briefly.

\begin{verbatim}
empty_dict = {}
empty_set = set()

print("empty_dict")
print(f"  type: {type(empty_dict)}")
print(f"  keys in dir: {'keys' in dir(empty_dict)}")
print(f"  items in dir: {'items' in dir(empty_dict)}")
print(f"  add in dir: {'add' in dir(empty_dict)}")
print()

print("empty_set")
print(f"  type: {type(empty_set)}")
print(f"  keys in dir: {'keys' in dir(empty_set)}")
print(f"  items in dir: {'items' in dir(empty_set)}")
print(f"  add in dir: {'add' in dir(empty_set)}")
\end{verbatim}

Possible output:

\begin{verbatim}
empty_dict
  type: <class 'dict'>
  keys in dir: True
  items in dir: True
  add in dir: False

empty_set
  type: <class 'set'>
  keys in dir: False
  items in dir: False
  add in dir: True
\end{verbatim}

The dictionary has \texttt{keys()} and \texttt{items()} because it stores key-value associations. The set has \texttt{add()} because it stores individual unique elements.

A single value inside curly braces creates a set, not a dictionary.

\begin{verbatim}
one_item = {"Serenity now!"}

print(f"one_item: {one_item!r}")
print(f"type(one_item): {type(one_item)}")
\end{verbatim}

Possible output:

\begin{verbatim}
one_item: {'Serenity now!'}
type(one_item): <class 'set'>
\end{verbatim}

A colon creates dictionary syntax.

\begin{verbatim}
one_pair = {"Jerry": "Serenity now!"}

print(f"one_pair: {one_pair!r}")
print(f"type(one_pair): {type(one_pair)}")
\end{verbatim}

Possible output:

\begin{verbatim}
one_pair: {'Jerry': 'Serenity now!'}
type(one_pair): <class 'dict'>
\end{verbatim}

The syntax boundary is therefore:

\begin{verbatim}
{"spam"}          -> set with one element
{"spam", "eggs"}  -> set with two elements
{}                -> empty dict
set()             -> empty set
{"name": "spam"}  -> dict with one key-value pair
\end{verbatim}

This boundary matters because set code and dictionary code support different operations. Confusing \texttt{\{\}} with an empty set often leads to errors later, when the program tries to call set-specific methods on a dictionary.

\begin{verbatim}
items = {}

print(f"type(items): {type(items)}")

try:
    items.add("spam")
except AttributeError as err:
    print(f"method call failed: {err}")
\end{verbatim}

Possible output:

\begin{verbatim}
type(items): <class 'dict'>
method call failed: 'dict' object has no attribute 'add'
\end{verbatim}

The fix is to use \texttt{set()} when the intended object is an empty set.

\begin{verbatim}
items = set()

items.add("spam")

print(f"items: {items!r}")
print(f"type(items): {type(items)}")
\end{verbatim}

Possible output:

\begin{verbatim}
items: {'spam'}
type(items): <class 'set'>
\end{verbatim}

The mutation operation \texttt{add()} will be studied later in more detail. Here it is used only to expose the syntactic difference between an empty dictionary and an empty set.

\subsubsection{Set Comprehensions as Hash-Based Filtering Structures}

A set comprehension builds a set by evaluating an expression across an iterable source. Its syntax resembles a list comprehension, but it uses curly braces and produces a set.

\begin{verbatim}
words = ["spam", "eggs", "spam", "larch", "eggs"]

unique_words = {word for word in words}

print(f"words: {words!r}")
print(f"type(words): {type(words)}")
print()

print(f"unique_words: {unique_words!r}")
print(f"type(unique_words): {type(unique_words)}")
print(f"len(unique_words): {len(unique_words)}")
\end{verbatim}

Possible output:

\begin{verbatim}
words: ['spam', 'eggs', 'spam', 'larch', 'eggs']
type(words): <class 'list'>

unique_words: {'eggs', 'spam', 'larch'}
type(unique_words): <class 'set'>
len(unique_words): 3
\end{verbatim}

The source list has five components. The resulting set has three elements because repeated values collapse into one set member.

A set comprehension can also include a filtering condition.

\begin{verbatim}
nums = [1, 2, 3, 4, 5, 42, 42]

large_nums = {num for num in nums if num > 3}

print(f"nums: {nums!r}")
print(f"large_nums: {large_nums!r}")
print(f"type(large_nums): {type(large_nums)}")
\end{verbatim}

Possible output:

\begin{verbatim}
nums: [1, 2, 3, 4, 5, 42, 42]
large_nums: {42, 4, 5}
type(large_nums): <class 'set'>
\end{verbatim}

This does two things at once:

\begin{itemize}
    \item it filters the input values with \texttt{if num > 3};
    \item it removes duplicates because the result is a set.
\end{itemize}

The same idea works with transformed values.

\begin{verbatim}
quotes = [
    "D'oh!",
    "d'oh!",
    "SPAM",
    "spam",
    "Nobody expects the Spanish Inquisition!",
]

normalized = {text.lower() for text in quotes}

print("original values")
for text in quotes:
    print(f"  {text!r}")

print()
print(f"normalized: {normalized!r}")
print(f"type(normalized): {type(normalized)}")
\end{verbatim}

Possible output:

\begin{verbatim}
original values
  "D'oh!"
  "d'oh!"
  'SPAM'
  'spam'
  'Nobody expects the Spanish Inquisition!'

normalized: {"d'oh!", 'spam', 'nobody expects the spanish inquisition!'}
type(normalized): <class 'set'>
\end{verbatim}

The transformation \texttt{text.lower()} runs before set insertion. After that, the hash-based uniqueness rule applies. Therefore, \texttt{"SPAM"} and \texttt{"spam"} become the same resulting value.

Set comprehensions are useful when the program needs a unique domain derived from another structure.

\begin{verbatim}
names = ["Homer", "Marge", "Lisa", "Homer", "Bart", "Lisa"]

name_lengths = {len(name) for name in names}

print(f"names: {names!r}")
print(f"name_lengths: {name_lengths!r}")
print(f"type(name_lengths): {type(name_lengths)}")
\end{verbatim}

Possible output:

\begin{verbatim}
names: ['Homer', 'Marge', 'Lisa', 'Homer', 'Bart', 'Lisa']
name_lengths: {4, 5}
type(name_lengths): <class 'set'>
\end{verbatim}

The result does not say which name had which length. It only says which distinct lengths occurred. That is the correct set-level meaning.

A set comprehension should be read as:

\begin{verbatim}
build a set
for each source element
    compute the result expression
    optionally reject it with an if condition
    insert the result into the hash-based uniqueness domain
\end{verbatim}

So the comprehension:

\begin{verbatim}
{word.lower() for word in words if len(word) > 4}
\end{verbatim}

means:

\begin{itemize}
    \item traverse \texttt{words};
    \item keep only words whose length is greater than \texttt{4};
    \item convert each kept word to lowercase;
    \item insert the lowercase result into a set;
    \item collapse duplicates automatically.
\end{itemize}

A final example shows this complete path.

\begin{verbatim}
words = [
    "Spam",
    "eggs",
    "SPAM",
    "larch",
    "LARCH",
    "Ni!",
]

result = {word.lower() for word in words if len(word) > 3}

print(f"words: {words!r}")
print(f"result: {result!r}")
print(f"type(result): {type(result)}")
print(f"len(result): {len(result)}")
\end{verbatim}

Possible output:

\begin{verbatim}
words: ['Spam', 'eggs', 'SPAM', 'larch', 'LARCH', 'Ni!']
result: {'eggs', 'spam', 'larch'}
type(result): <class 'set'>
len(result): 3
\end{verbatim}

The result is not a filtered list. It is a filtered, transformed, hash-based uniqueness domain.